\section{Weighted reflection-positive tuple laws}
\label{sec:tuples}

We now construct the finite-field probability laws used to average products
of Fourier lifts.  The two requirements that must be reconciled are
positivity across certain reflections of the slots and support on
polynomial differences around a prescribed cycle.  A measure obtained
simply by imposing the cycle constraints need not have the required
positivity.  We instead construct several measures, indexed by sets of
\emph{marked blocks}.  Marking a block removes some constraints and creates
a positive contribution at reflections in that block.  A suitably weighted
mixture lets this contribution dominate the possible negative contribution
from the corresponding unmarked measure.

The graph estimates and marked-block construction adapt
\cite[Lemmas~2.1--2.2 and Section~3]{squarepower}.  Here the edge constraints
are weighted by regular fibers of $h_l$.  Their directions are retained
throughout: neither the polynomial value set nor its weight is assumed to
be invariant under negation.  We prove all the properties needed below,
including the exact covariance under changes of the auxiliary parameter.

\subsection{Slots, reflections, and mixing of edge weights}

Retain the constants $\delta,\gamma$ and the integer $L>2$ from
Section~\ref{sec:local}.  Choose the least integer $t\ge1$ such that
$(t+1)\delta/2>5$, and set
\[
 w=2t+1,\qquad K=L!,\qquad r_*=K/2,\qquad
 V=\Pi^w,\qquad d=K^w,\qquad n_*=d/2,
\]
where $\Pi$ is the set of bijections from the positions $[L]$ to the
symbols $[L]$.  Thus an element of $\Pi$ is a permutation word, and a slot
$v\in V$ consists of $w$ such words, which we call its blocks.  Fix an
$L$-cycle $c$ acting on positions.  Right composition $q\mapsto qc$
shifts a word, whereas left composition permutes its symbols.  The
designated cycle of slots is
\[
 v_j=(c^j,\ldots,c^j)\quad(0\le j<L),\qquad v_L=v_0.
\]

A reflection cut is specified by a block $b\in[w]$ and two distinct
symbols $a_1,a_2$ in a specified order.  Let $V_+$ consist of the slots
in which $a_1$ precedes $a_2$ in block $b$.  Let $\theta$ interchange
these two symbols in that block, leaving every other block unchanged.
Then $\theta$ is an involution exchanging $V_+$ and
$V_-=V\setminus V_+$, and both halves have size $n_*$.  For a finite
measure $\nu$ on $\F_p^V$, its matrix at this cut is the matrix of joint
point masses of
\[
 (z_v)_{v\in V_+}
 \quad\hbox{and}\quad
 (z_{\theta v})_{v\in V_+}.
\]
We call $\nu$ \emph{reflection positive} if this matrix is symmetric
positive semidefinite at every cut.  Equivalently, symmetry holds and
\[
 \int J((z_v)_{v\in V_+})
          J((z_{\theta v})_{v\in V_+})\,d\nu(z)\ge0
\]
for every real function $J$ on $\F_p^{V_+}$ and every cut.

Fix $l\ge1$ and a prime $p\ge P_1$, and write $F=h_l$.  Define
\[
 W(y)=W_{l,p}(y)=\frac1k
   \bigl|\{x\in\F_p:F(x)=y,\ F'(x)\ne0\}\bigr|,
 \qquad c_p=\E_{y\in\F_p}W(y).
\]
The mean $c_p$ may depend on $l$.  Nonconstancy of $F$ modulo $p$ and
the degree bound give $0\le W\le1$, and $W$ is supported on the
admissible residues.  Moreover,
\[
 c_p=\frac{\bigl|\{x:F'(x)\ne0\}\bigr|}{kp}
       \ge\frac1{2k}.
\]
Equation~\eqref{eq:2} bounds every nonconstant Fourier coefficient
of $W$ by $p^{-4\delta}$.  Consequently the kernel
$W(y-x)-c_p$, acting between uniform $L^2(\F_p)$ spaces, has operator
norm at most $p^{-4\delta}$.  This follows by diagonalizing the
convolution operator in the character basis; no symmetry of $W$ is
needed.

We shall apply this one-edge estimate to fixed systems of edges.  All
graphs in the next lemma have distinct formal vertices, whose field
labels are independent and uniform before the edge weights are applied.

\begin{lemma}[Graph counting and list mixing]
\label{tuple:list-mixing}
Let $G$ be a fixed directed graph with no loops and with at most one
edge on each unordered pair of vertices.  Write $e_G$ for its number
of edges.  All averages below use the normalized uniform measures on
the indicated label configurations.  Then
\[
 \left|\E\prod_{x\to y\in E(G)}W(Y_y-Y_x)-c_p^{e_G}\right|
       \le e_Gp^{-4\delta}.
\]
If every edge joins a vertex in one list $X$ to a vertex in a disjoint
list $Y$, in either direction, let $K_G(X,Y)$ denote this product of
edge weights.  As a kernel acting by normalized averaging between the
uniform $L^2$ spaces of the two lists,
\begin{equation}
 \|K_G-c_p^{e_G}\|_{\mathrm{op}}
       \le(32e_G)^{1/4}p^{-\delta}.
 \label{eq:3}
\end{equation}
Both bounds are uniform in $l$ and $p$ and include $e_G=0$, when the
products are one and both errors are zero.
\end{lemma}

\begin{proof}
Reversing a directed edge replaces the one-edge operator by its adjoint,
so its norm remains at most $p^{-4\delta}$.  The case $e_G=0$ is
immediate; assume $e_G\ge1$.
For the first assertion, replace the edge weights by $c_p$ one at a
time.  In a term containing the difference for an edge $x\to y$, fix
all labels other than $Y_x,Y_y$.  Because there is no other edge on
this unordered pair, the remaining factors separate into a function
of $Y_x$ and a function of $Y_y$, each bounded by one and hence of
uniform $L^2$ norm at most one.  The one-edge operator estimate bounds
this term by $p^{-4\delta}$.  Summing over the $e_G$ replacements gives
the stated error $e_Gp^{-4\delta}$.

For the second assertion put $B=K_G-c_p^{e_G}$.  If the two configuration
spaces have cardinalities $N_X,N_Y$, the kernel operator
$(T_Bf)(X)=\E_Y B(X,Y)f(Y)$ has Euclidean matrix
$B(X,Y)/\sqrt{N_XN_Y}$ under the isometries from uniform $L^2$ to
Euclidean space.  Thus, with two independent copies of each list,
matrix multiplication gives
\[
 \|B\|_{\mathrm{op}}^4
 \le \sum_j s_j(B)^4
 =\E\prod_{a,b\in\{0,1\}}B(X_a,Y_b),
\]
where $s_j(B)$ are the singular values for the normalized measures.
Indeed, the trace of $(T_BT_B^*)^2$ is the uniform average of
$B(X_0,Y_0)\overline{B(X_1,Y_0)}B(X_1,Y_1)\overline{B(X_0,Y_1)}$;
here $B$ is real, so the conjugates can be omitted.  This explains why
no factor depending on $N_X,N_Y$ occurs.
Expand the four copies of $B$.  A product using any subset of the four
corners again comes from a graph of the stated kind: an unordered pair
of its formal endpoints determines both the original edge and its
corner.  A term using $j$ corners has exactly $je_G$ edges, so its graph
error is at most $je_Gp^{-4\delta}$.  Its constant multiplier has
modulus $c_p^{(4-j)e_G}\le1$.  All main terms are $c_p^{4e_G}$ times
the corresponding sign, so they cancel.  Consequently
\[
 \|B\|_{\mathrm{op}}^4
 \le\sum_{j=0}^4\binom4j je_Gp^{-4\delta}
 =32e_Gp^{-4\delta},
\]
because $\sum_{j=0}^4\binom4j j=4\cdot2^3=32$.  Taking fourth roots
proves \eqref{eq:3}.
\end{proof}

\subsection{Measures indexed by marked blocks}

For $S\subseteq[w]$ with $s=|S|\le t+1$, regard the blocks in $S$ as
marked.  The \emph{address} of a slot is the tuple of its unmarked
words:
\[
 u_S:V\longrightarrow U_S=\Pi^{[w]\setminus S}.
\]
At each address take a list of length $\ell_s=r_*^s$,
\[
 (Y_{u,i})_{i\in[\ell_s]},\qquad u\in U_S,
\]
with all entries independent and uniform in $\F_p$.  The length grows
with $s$ so that, upon marking one more block, the lists on one side
of a cut can be concatenated into the new lists.

For $s\le t$, put a directed edge $u\to u'$ between addresses whenever
\[
 u'_b=u_bc
 \quad\hbox{in at least $t+1$ blocks }b\notin S.
\]
For every address edge impose all list-entry weights in its direction,
and write their product as
\[
 I_S(Y)=
 \prod_{u\to u'}\ \prod_{i,i'\in[\ell_s]}
        W(Y_{u',i'}-Y_{u,i}).
\]
There are no loops, since $qc\ne q$ for a permutation word $q$.
Nor can both $u\to u'$ and $u'\to u$ occur: a block certifying one
direction cannot certify the other, as $c^2\ne1$.  The two directions
would therefore require at least $2(t+1)>w-s$ unmarked blocks.
The resulting graph on list entries satisfies
Lemma~\ref{tuple:list-mixing}.  The same threshold also preserves the
designated cycle: if $s\le t$, at least $t+1$ blocks remain unmarked,
and each consecutive pair of cycle vertices shifts in every one of
them.  At $s=t+1$ only $t$ blocks remain, and we set $I_S=1$;
we call these sets $S$ terminal.

Each slot will choose an entry of its address list.  To obtain strict
positivity at a marked cut, we also weight the choices themselves.
Set $\Delta=1/(10L^2)$.  For an unordered pair of distinct slots
$x,y\in V$, set $H^S_{xy}=0$ unless $x,y$ agree in every block except
one marked block and their words in that block differ by a
transposition of symbols.  If the exchanged positions have gap $g$,
set $H^S_{xy}=\Delta^g$.  Choose independent uniform indices
$j_v\in[\ell_s]$ and put
\[
 R_S(j)=\exp\left(\sum_{\{x,y\}\subseteq V}
                         H^S_{xy}\1_{\{j_x=j_y\}}\right).
\]
Define a finite, not yet normalized, measure $\nu_S=\nu_{l,p,S}$ by
\[
 \int \Phi(z)\,d\nu_S(z)
   =\E_{Y,j}\left[
       I_S(Y)R_S(j)
       \Phi\bigl((Y_{u_S(v),j_v})_{v\in V}\bigr)\right].
\]
Thus the output label at $v$ is $z_v=Y_{u_S(v),j_v}$.  All underlying
averages in this definition are independent uniform averages.
Since the index set is fixed,
\[
 0\le I_S\le1,\qquad 1\le R_S\le C,
 \qquad \nu_S(1)\le C,
\]
with a constant independent of $l,p,S$.

Every $\nu_S$ is invariant under the slot involution $\theta$ of any
reflection cut.  Indeed, simultaneous symbol relabeling preserves
$H^S$: it conjugates the transposition defining an interaction and
leaves the exchanged positions unchanged.  If the cut block is
marked, addresses are unchanged.  If it is unmarked, reflection
permutes the lists, preserving each directed edge because left symbol
permutations commute with right composition by $c$.  Reindexing the
independent lists and indices therefore preserves the measure.  In
particular, every component has a symmetric matrix at every cut.

\subsection{Positivity across a cut}

Fix a reflection cut in block $b$.  For a real function
$J:\F_p^{V_+}\to\R$, write
\[
 Q_S(J)=\int J((z_v)_{v\in V_+})
                 J((z_{\theta v})_{v\in V_+})\,d\nu_S(z).
\]
We compare $Q_S$ with $Q_{S\cup\{b\}}$.  When $b$ is marked,
the index interactions give a fixed positive lower bound.  When $b$
is unmarked, list mixing bounds any negative contribution by a power
of $p$.  The same half-label square occurs in both bounds after
concatenating the lists.

\begin{lemma}[Positivity at a marked block]
\label{tuple:marked-positive}
There is a constant $c_*>0$, independent of $l,p,S$ and the cut, such
that if $b\in S$, then
\begin{equation}
 Q_S(J)\ge c_*\,
 \E_{Y,j_+} I_S(Y)
       J\bigl((Y_{u_S(v),j_v})_{v\in V_+}\bigr)^2.
 \label{eq:4}
\end{equation}
Here $j_+$ consists of independent uniform indices on $V_+$.
\end{lemma}

\begin{proof}
The interaction between the two halves is governed by the symmetric
matrix
\[
 \mathsf A_{xy}=H^S_{x,\theta y},\qquad x,y\in V_+.
\]
Fix $x\in V_+$, and suppose that the cut symbols occur at positions
$i<j$ in its cut block, with gap $g=j-i$.  Swapping the cut symbols
gives the diagonal entry $\mathsf A_{xx}=\Delta^g$.  Any other
transposition carrying $x$ across the cut must move the first cut
symbol to a position beyond $j$, or the second to a position before
$i$.  Its gap is at least $g+1$.  There are fewer than $L^2$ such
transpositions; transpositions in other blocks do not cross the cut.
Consequently the absolute off-diagonal row sum is at most
$L^2\Delta^{g+1}=\Delta^g/10$.  Symmetric diagonal dominance gives
\[
 \mathsf A\succeq\kappa\operatorname{Id},
 \qquad \kappa=\frac9{10}\Delta^{L-1}>0.
\]

For $j\in[\ell_s]^{V_+}$ define
$u(j)=(\1_{\{j_x=i\}})_{x\in V_+,\,i\in[\ell_s]}$.
After identifying the minus side with the plus side by reflection,
the cross-interaction matrix on half-index assignments is
\[
 C_{j,j'}=\exp\bigl(u(j)^T
              (\mathsf A\otimes\operatorname{Id})u(j')\bigr).
\]
Split $\mathsf A=(\mathsf A-\kappa\operatorname{Id})+
\kappa\operatorname{Id}$.  The entrywise exponential associated with
the first summand is positive semidefinite: its power-series
expansion is a sum of Gram matrices of tensor powers.  Its diagonal
entries are at least one.  The exponential associated with the second
summand is a tensor product, over $x\in V_+$, of
\[
 \mathbf1\mathbf1^T+(\exp(\kappa)-1)\operatorname{Id}
\]
on $[\ell_s]$.  It therefore dominates
$a_*\operatorname{Id}$, where
$a_*=(\exp(\kappa)-1)^{n_*}$.  The entrywise product of positive
semidefinite matrices is positive semidefinite, as follows directly
from their Gram representations.  Multiply the lower bound for the
second exponential entrywise by the first.  The resulting diagonal
lower bound is at least $a_*\operatorname{Id}$ because the first
exponential has diagonal entries at least one.  Thus
$C\succeq a_*\operatorname{Id}$.

The interactions within the two halves give the same diagonal factor
on either side of $C$, with all entries at least one.  The full index
matrix retains the lower bound $a_*\operatorname{Id}$.  For fixed
lists $Y$, reflected vertices have the same address because $b$ is
marked.  The half-label evaluations therefore form the same vector
on both sides of this matrix, even when some field labels coincide.
There are $\ell_s^{n_*}$ half-index assignments.  Averaging uniformly
on both sides converts the matrix lower bound into
\[
 \E_{j_+,j_-}R_S(j)
       J((z_v)_{v\in V_+})J((z_{\theta v})_{v\in V_+})
 \ge a_*\ell_s^{-n_*}\,
       \E_{j_+}J((z_v)_{v\in V_+})^2.
\]
Multiplying by $I_S(Y)$ and averaging $Y$ proves \eqref{eq:4} with
$c_*=a_*\ell_{t+1}^{-n_*}$.
\end{proof}

\begin{lemma}[The contribution at an unmarked block]
\label{tuple:unmarked-defect}
Suppose $b\notin S$.  Split the address lists into the two halves of
the cut.  Denote the plus lists by $Y_+$, and let $I_+(Y_+)$ be the
product of the weights on edges internal to them.  Put
\[
 \mathcal D_S(J)=\E_{Y_+,j_+} I_+(Y_+)
       J\bigl((Y_{u_S(v),j_v})_{v\in V_+}\bigr)^2.
\]
If $|S|\le t$, then
\begin{equation}
 Q_S(J)\ge-Cp^{-\delta}\mathcal D_S(J).
 \label{eq:5}
\end{equation}
If $|S|=t+1$, then $Q_S(J)\ge0$.
\end{lemma}

\begin{proof}
Reflection identifies the minus lists and their internal edge factors
with the plus lists and their internal factors.  No index interaction
crosses this cut: an interaction changes only a marked block, whereas
the cut is determined by the unmarked block $b$.  The two within-half
index weights are likewise identified by reflection; write $R_+$
for the one on the plus half.

Let $K(Y_+,Y_-)$ be the product of all cross-list edge weights.  Then
$Q_S(J)$ is the quadratic form of this kernel on the half function
\[
 g(Y_+)=I_+(Y_+)\E_{j_+}
       R_+(j_+)J\bigl((Y_{u_S(v),j_v})_{v\in V_+}\bigr).
\]
Cauchy--Schwarz, the uniform bound on $R_+$, and $I_+^2\le I_+$ give
$\|g\|_2^2\le C\mathcal D_S(J)$.  By \eqref{eq:3}, the kernel
$K$ differs in operator norm by $O(p^{-\delta})$ from the constant
$c_p^{e_{\rm cross}}$, where $e_{\rm cross}$ is its number of edges.
The constant contributes $c_p^{e_{\rm cross}}(\E g)^2\ge0$.
This proves \eqref{eq:5}.  In the terminal case $I_S=1$, so the
cross kernel is exactly one and the entire form is nonnegative.
\end{proof}

\begin{lemma}[Domination after marking one more block]
\label{tuple:marking-domination}
If $b\notin S$ and $|S|\le t$, then, with $S'=S\cup\{b\}$,
\begin{equation}
 Q_{S'}(J)\ge c_*r_*^{-n_*}\mathcal D_S(J).
 \label{eq:6}
\end{equation}
\end{lemma}

\begin{proof}
Let $\Pi_+$ be the $r_*$ words in block $b$ satisfying the cut order.
For each new address $u\in U_{S'}$, the old plus addresses above it
are $(u,q)$ with $q\in\Pi_+$; here $q$ is inserted in block $b$.
Concatenate their lists by identifying
\[
 [\ell_{s+1}]\ \cong\ \Pi_+\times[\ell_s],
 \qquad
 Y'_{u,(q,i)}=Y_{(u,q),i}.
\]
This is a bijection of independent uniform list entries, so it changes
neither their joint distribution nor the value of an average.
Figure~\ref{tuple:fig-regrouping} separates this relabeling from the
index restriction used below.

\begin{figure}[htbp]
\centering
\begin{tikzpicture}[
  x=1cm,y=1cm,
  font=\small,
  list/.style={draw=black!65,rounded corners=1pt,
               minimum width=3.1cm,minimum height=.56cm},
  selected/.style={fill=blue!9,draw=blue!60!black},
  every path/.style={line width=.55pt}
]
  \node[align=center] at (1.55,.85)
    {$r_*$ old plus lists\\addresses $(u,q)$, $q\in\Pi_+$};
  \node[align=center] at (8.65,.85)
    {one new list at $u$\\length $\ell_{s+1}=r_*\ell_s$};

  \node[list] at (1.55,0) {$1\quad\cdots\quad\ell_s$};
  \node[anchor=east] at (-.15,0) {$q_1$};
  \node at (1.55,-.62) {$\vdots$};
  \node[list,selected] at (1.55,-1.25)
    {$1\quad\cdots\quad\ell_s$};
  \node[anchor=east] at (-.15,-1.25) {$q=v_b$};
  \node at (1.55,-1.88) {$\vdots$};
  \node[list] at (1.55,-2.5) {$1\quad\cdots\quad\ell_s$};
  \node[anchor=east] at (-.15,-2.5) {$q_{r_*}$};

  \draw[<->,>=Stealth] (3.38,-1.25)--(6.82,-1.25);
  \node[above,align=center] at (5.1,-1.13) {bijective relabeling};
  \node[below,align=center] at (5.1,-1.37) {no restriction};

  \draw[black!65,rounded corners=1pt] (7.05,.32) rectangle (10.25,-2.82);
  \node at (8.65,0) {$\{q_1\}\times[\ell_s]$};
  \node at (8.65,-.62) {$\vdots$};
  \filldraw[selected] (7.05,-.97) rectangle (10.25,-1.53);
  \node at (8.65,-1.25) {$\{v_b\}\times[\ell_s]$};
  \node at (8.65,-1.88) {$\vdots$};
  \node at (8.65,-2.5) {$\{q_{r_*}\}\times[\ell_s]$};

  \draw[->,>=Stealth] (8.65,-2.92)--(8.65,-3.72);
  \node[anchor=east,align=right] at (8.45,-3.32)
    {restrict the index\\to $q=v_b$};
  \node[anchor=west] at (8.85,-3.32) {probability $1/r_*$};
  \node[list,selected] at (8.65,-4.1) {$j'_v=(v_b,i)$};
  \node[below] at (8.65,-4.42) {$i$ uniform on $[\ell_s]$};
\end{tikzpicture}
\caption{Regrouping at a fixed new address
$u=u_{S\cup\{b\}}(v)$ and restricting the index of one slot $v\in V_+$.
The horizontal arrow only relabels the same independent entries; it
incurs no probability factor.  The highlighted sublist is prescribed
by $v_b$.  Restricting all $n_*$ independent slot indices to their
respective prescribed sublists has probability $r_*^{-n_*}$.
The ellipses denote the other words in $\Pi_+$; boxes describe index
positions, and their field values may coincide.}
\label{tuple:fig-regrouping}
\end{figure}


A new address edge requires $t+1$ shifted blocks outside $S'$.  Those blocks
also certify an old edge for every choice of the two words in block
$b$.  Hence every factor of $I_{S'}(Y')$ occurs, with its direction
unchanged, among the factors of $I_+(Y_+)$.  Since all factors belong
to $[0,1]$,
\[
 I_+(Y_+)\le I_{S'}(Y').
\]
When $S'$ is terminal this inequality follows from $I_{S'}=1$.

Apply \eqref{eq:4} to $S'$.  In its nonnegative right-hand side,
restrict the new index at each $v\in V_+$ to the sublist
$\{v_b\}\times[\ell_s]$.  Each restriction has probability $1/r_*$,
independently across the $n_*$ slots.  On this event the labels and
the remaining independent indices are exactly those in
$\mathcal D_S(J)$.  Its probability is $r_*^{-n_*}$, and the
pointwise weight comparison above now yields \eqref{eq:6}.
\end{proof}

\subsection{The mixture, its support, and covariance}

We can now choose the mixture.  The gain in \eqref{eq:6} is a fixed
positive constant, whereas the possible loss in \eqref{eq:5} tends
to zero with $p$.  Thus a component with one more marked block may
receive a smaller coefficient and still compensate for that loss.

\begin{proposition}[Reflection-positive laws with polynomial cycle support]
\label{tuple:laws}
For all sufficiently large primes $p$, uniformly in $l\ge1$, define
\begin{align*}
 M_{l,p}&=\sum_{\substack{S\subseteq[w]\\|S|\le t+1}}
                 p^{-|S|\delta/2}\nu_{l,p,S},\\
 \mu_{l,p}&=\frac{M_{l,p}}{M_{l,p}(1)},\\
 \eps_{l,p}&=\frac{1}{M_{l,p}(1)}
       \sum_{|S|=t+1}p^{-(t+1)\delta/2}\nu_{l,p,S}.
\end{align*}
Then $\mu_{l,p}$ is a reflection-positive probability law and
$0\le\eps_{l,p}\le\mu_{l,p}$.  Its exceptional mass satisfies
\[
 \eta_{l,p}:=\eps_{l,p}(1)\le p^{-4}.
\]
The probability law
\[
 \mu_{l,p}^{\circ}
       =\frac{\mu_{l,p}-\eps_{l,p}}{1-\eta_{l,p}}
\]
has only admissible increments $z_{v_{j+1}}-z_{v_j}$ along the
designated cycle.  Every component and both measures $\mu_{l,p}$,
$\eps_{l,p}$ are invariant under simultaneous translation of all
labels.  Every positive-mass component or subsum therefore has uniform
single-slot marginals after normalization.  In either $\mu_{l,p}$ or
$\mu_{l,p}^{\circ}$, the total probability of components with
$S\ne\varnothing$ is $O(p^{-\delta/2})$.
\end{proposition}

\begin{proof}
At a cut in block $b$, pair each $S$ with $b\notin S$, $|S|\le t$,
with $S\cup\{b\}$.  Each component containing $b$ appears in exactly
one pair.  Equations~\eqref{eq:5} and \eqref{eq:6} bound the quadratic
form of the pair below by
\[
 p^{-|S|\delta/2}
 \left(p^{-\delta/2}c_*r_*^{-n_*}-Cp^{-\delta}\right)
 \mathcal D_S(J).
\]
This is nonnegative once $p$ is sufficiently large.  The only unpaired
sets are terminal sets omitting $b$; their forms are nonnegative by
Lemma~\ref{tuple:unmarked-defect}.  Thus $M_{l,p}$ is reflection
positive at every cut.

For $S=\varnothing$, each list has length one and $R_S=1$.
Lemma~\ref{tuple:list-mixing} and $c_p\ge1/(2k)$ show that
$\nu_{\varnothing}(1)$ is bounded below by a fixed positive constant
for all sufficiently large $p$.  All component masses are bounded
above, and the number of components is fixed.  It follows that
\[
 0<c\le M_{l,p}(1)\le C,\qquad
 \eps_{l,p}(1)\le C p^{-(t+1)\delta/2}.
\]
Our choice $(t+1)\delta/2>5$ gives the stated bound $p^{-4}$ after
increasing the prime threshold.  Summing the nonleading coefficients
also gives total probability $O(p^{-\delta/2})$ for $S\ne\varnothing$.
Removing the terminal components retains the positive leading mass,
so the same estimate holds after normalization to $\mu_{l,p}^{\circ}$.

Simultaneous translation of all list entries preserves their product
law, all differences in $I_S$, and the index weight $R_S$.  It
therefore translates every output label and proves the asserted
invariances and marginal uniformity.

Finally, a nonterminal $S$ leaves at least $t+1$ blocks unmarked.
Consecutive cycle vertices shift by $c$ in all these blocks, including
the edge from $v_{L-1}$ to $v_0$.  Their addresses are consequently
joined by a directed edge.  Since its weights include every pair of
list entries, any positive-weight label tuple has admissible
increment on this edge, regardless of the chosen indices.  This
holds around the entire cycle in every component retained by
$\mu_{l,p}^{\circ}$.
\end{proof}

The auxiliary parameter changes the edge weight, so we need covariance
between different laws rather than dilation invariance of a single law.

\begin{lemma}[Covariance under progression changes]
\label{tuple:covariance}
For every $l\ge1$, every sufficiently large prime $p$, and every
integer $D\ge1$ with $p\nmid D$, the following holds.  For each
$a'\in\F_p$, the map
\begin{equation}
 (z_v)_{v\in V}\longmapsto
       \bigl(\lambda(D)^{-1}(z_v-a')\bigr)_{v\in V}
 \label{eq:7}
\end{equation}
pushes $\mu_{l,p}$ to $\mu_{lD,p}$ and $\eps_{l,p}$ to
$\eps_{lD,p}$.
\end{lemma}

\begin{proof}
Equation~\eqref{eq:1} states that
$h_l(Dx-i)=\lambda(D)h_{lD}(x)$.  Both $D$ and $\lambda(D)$ are
units modulo $p$.  The argument map $x\mapsto Dx-i$ is therefore a
bijection and, on differentiating this identity, preserves regularity.
Counting regular arguments gives
\[
 W_{lD,p}(y)=W_{l,p}(\lambda(D)y).
\]
Apply \eqref{eq:7} to every list entry in a component measure.  The
independent uniform list law and all index weights are unchanged;
each old edge weight becomes exactly the corresponding new edge
weight.  Thus the map pushes $\nu_{l,p,S}$ to $\nu_{lD,p,S}$ for
each $S$.  All component masses, and in particular the normalizing
masses $M_{l,p}(1)$ and $M_{lD,p}(1)$, agree.  The two claimed
pushforwards follow with no additional factor.
\end{proof}

\subsection{Conditional mixing and the prime threshold}

The mixture has the required positivity and support.  To control high
Fourier frequencies later, we also need to know that the other labels
give little information about a mean-zero function of any one label.
The leading component has this property by list mixing.  The remaining
components have small total probability, and their uniform marginals
give the bound needed to include them.

\begin{lemma}[Conditional mixing]
\label{tuple:conditional-mixing}
For all sufficiently large $p$, uniformly in $l$, let $\rho_p$ be
either $\mu_{l,p}$ or $\mu_{l,p}^{\circ}$.  For every $v\in V$ and
every function $f:\F_p\to\C$ of uniform mean zero,
\begin{equation}
 \left\|\E_{\rho_p}\bigl(f(z_v)\mid
                         (z_{v'})_{v'\ne v}\bigr)\right\|_2
       \le p^{-\gamma}\|f\|_2.
 \label{eq:8}
\end{equation}
The norm on the left is taken in the other-label marginal of
$\rho_p$; the norm on the right is uniform on $\F_p$.
\end{lemma}

\begin{proof}
First normalize the leading component $\nu_{\varnothing}$, and
write $x=z_v$, $y=(z_{v'})_{v'\ne v}$.  Before normalization its
density relative to uniform labels is $K_1(x,y)K_2(y)$, where $K_1$
is the product of the $m$ weights on edges incident to $v$, and
$K_2$ contains the remaining weights.  Put
\[
 A(y)=\E_xK_1(x,y),\qquad
 B_f(y)=\E_x f(x)K_1(x,y).
\]
The incident edges form a graph between the singleton list $x$ and
the list $y$.  Equation~\eqref{eq:3}, applied to this kernel, gives
\[
 \|A-c_p^m\|_2\ll p^{-\delta},\qquad
 \|B_f\|_2\ll p^{-\delta}\|f\|_2.
\]
This includes $m=0$, when $A=1$ and $B_f=0$.

Let $Z_0=\nu_{\varnothing}(1)$, which is bounded below uniformly.
The other-label marginal has density $Z_0^{-1}K_2(y)A(y)$, and the
conditional expectation equals $B_f(y)/A(y)$ where $A(y)>0$.
Its squared norm is therefore exactly
\[
 \frac1{Z_0}\E_y K_2(y)\frac{|B_f(y)|^2}{A(y)},
\]
with the ratio taken as zero if $A(y)=0$.  On $A\ge c_p^m/2$, the
bound on $B_f$ makes this $O(p^{-2\delta})\|f\|_2^2$.  The
complement has uniform probability $O(p^{-2\delta})$, by the bound
on $A-c_p^m$ and the fixed positive lower bound on $c_p^m$.
Everywhere, weighted Cauchy--Schwarz gives
\[
 \frac{|B_f(y)|^2}{A(y)}
       \le\E_x |f(x)|^2K_1(x,y)\le\|f\|_2^2.
\]
Since $K_2\le1$, the complement has the same required bound.  Thus
the squared conditional norm in the normalized leading component is
$O(p^{-2\delta})\|f\|_2^2$.

For either mixture, condition additionally on its component index.
The tower property and Jensen's inequality can only increase the
squared conditional norm under this extra conditioning.  In each
nonleading component the trivial bound is $\|f\|_2^2$, because its
single-slot marginal is uniform.  Proposition~\ref{tuple:laws} bounds
their total probability by $O(p^{-\delta/2})$.  The squared norm for
the mixture is consequently $O(p^{-\delta/2})\|f\|_2^2$, and the
norm is $O(p^{-\delta/4})\|f\|_2$.  Since $\gamma=\delta/10$, the
fixed implicit constant is absorbed by taking $p$ sufficiently large,
proving \eqref{eq:8}.
\end{proof}

We require a threshold $P>P_1$ such that all the preceding
assertions hold for every prime $p\ge P$, uniformly in $l$.  We
further require
\begin{equation}
 3\le p^{\gamma/2}\quad(p\ge P),\qquad
 \prod_{\substack{p\ge P\\p\ \mathrm{prime}}}(1+p^{-3})-1
       \le\frac14.
 \label{eq:9}
\end{equation}
The second requirement is possible by convergence of the product.
For such a choice, define the small-prime modulus
\[
 M_0=\prod_{\substack{p<P\\p\ \mathrm{prime}}}p^{E_p}\ge2,
\]
using the exponents $E_p$ from Section~\ref{sec:local}.  If a modulus
is $M_0$ times a squarefree product of primes at least $P$, a residue
modulo that modulus is called admissible for $l$ when it is admissible
at each of its local prime-power factors.  This convention combines
the small-prime lifting conditions with the cycle support supplied
by Proposition~\ref{tuple:laws}.

\subsection{Dependence of the tuple threshold}

To obtain an exponent depending only on the degree, we need the tuple
threshold to be uniform when the local choices are uniform.  The next
proposition traces that dependence while retaining the construction for
general fixed polynomial data.

\begin{proposition}[Dependence of the prime threshold]
\label{tuple:threshold-dependence}
For the auxiliary family and local choices above, the threshold $P$ can
be chosen as a function of $k,P_1,L,t$ alone so that
Proposition~\ref{tuple:laws}, Lemmas~\ref{tuple:covariance}
and~\ref{tuple:conditional-mixing}, and \eqref{eq:9} hold uniformly in
$l$.  For general fixed data, $P_1$ may depend on $h$, and $L$ may depend
on $h$ and its chosen roots.  If $h$ satisfies Condition~\eqref{local:unit-condition}
and the degree-dependent choices of $P_1$, $E_p$, and $L$ from
Section~\ref{sec:local} are used, then $P$ and $M_0$ depend only on $k$.
\end{proposition}

\begin{proof}
For every prime $p\ge P_1$, the quantitative local input to the estimates
above is
\[
 0\le W_{l,p}\le1,\qquad c_p\ge(2k)^{-1},\qquad
 \max_{u\ne0}|\widehat W_{l,p}(u)|\le p^{-4\delta}.
\]
Indeed, with the normalized Fourier convention, for $u\ne0$ the coefficient
$\widehat W_{l,p}(u)$ is $c_p$ times the corresponding regular-argument
average in \eqref{eq:2}; also $c_p\le1$.  These statements hold uniformly
in $l$.  We extract
constants from the preceding proofs using
\[
 e_*:=\binom d2,\qquad R_*:=\exp(e_*),\qquad
 C_{\mathrm{gr}}:=(32e_*)^{1/4},\qquad
 Z_*:=\frac12(2k)^{-e_*}.
\]
They depend only on $k,L,t$, since $d=(L!)^{2t+1}$.

For a set of $s$ marked blocks, the graph on list entries has
$K^{w-s}\ell_s=d/2^s\le d$ formal vertices.  Its simplicity was checked
above before field labels were assigned, so coincidences between their
values do not affect the graph bound.  Every graph for $I_S$, every
cross-list graph in Lemma~\ref{tuple:unmarked-defect}, and every incident
graph in Lemma~\ref{tuple:conditional-mixing} consequently has at most
$e_*$ edges.  Lemma~\ref{tuple:list-mixing} therefore bounds their graph
errors by $e_*p^{-4\delta}$.  For the cross-list and incident graphs,
which join two lists, it bounds the operator errors by
$C_{\mathrm{gr}}p^{-\delta}$.

Each index interaction $H^S_{xy}$ is at most one, so $R_S$ and every
within-half weight $R_+$ are at most $R_*$.  In particular,
$\nu_S(1)\le R_*$.  The proof of Lemma~\ref{tuple:marked-positive} gives
the explicit constant
\[
 c_*=(\exp(\kappa)-1)^{n_*}\ell_{t+1}^{-n_*},
 \qquad
 \kappa=\frac9{10}(10L^2)^{-(L-1)}.
\]
The exact half-index factor $\ell_s^{-n_*}$ proved there and the exact
index-restriction factor $r_*^{-n_*}$ in
Lemma~\ref{tuple:marking-domination} give the domination constant
$c_*r_*^{-n_*}$.  In the unmarked
proof, Cauchy--Schwarz and $I_+^2\le I_+$ give
$\|g\|_2^2\le R_*^2\mathcal D_S(J)$.  The list bound just obtained thus
allows the constant in \eqref{eq:5} to be
\[
 C_{\mathrm{def}}:=R_*^2C_{\mathrm{gr}}.
\]
The pairing proof of Proposition~\ref{tuple:laws} is therefore
nonnegative at every cut whenever
\[
 p^{\delta/2}\ge C_{\mathrm{def}}r_*^{n_*}/c_*.
\]

Let $e_0\le e_*$ be the number of edges in the leading component.
The graph estimate and the lower bound on $c_p$ give
\[
 \nu_{\varnothing}(1)
 \ge c_p^{e_0}-e_0p^{-4\delta}
 \ge (2k)^{-e_*}-e_*p^{-4\delta}.
\]
Consequently $\nu_{\varnothing}(1)\ge Z_*$ once
$e_*p^{-4\delta}\le Z_*$.  There are at most $2^w$ components, each of
mass at most $R_*$.  Set
\[
 C_{\mathrm{mass}}:=\frac{2^wR_*}{Z_*},\qquad
 \alpha_{\mathrm{exc}}:=\frac{(t+1)\delta}{2}>5.
\]
The normalizing mass $M_{l,p}(1)$ is at least the leading mass, so
\[
 \eta_{l,p}\le C_{\mathrm{mass}}p^{-\alpha_{\mathrm{exc}}}.
\]
It is at most $p^{-4}$ if
$p^{\alpha_{\mathrm{exc}}-4}\ge C_{\mathrm{mass}}$.  In either
$\mu_{l,p}$ or $\mu_{l,p}^{\circ}$ the probability of nonleading
components is at most $C_{\mathrm{mass}}p^{-\delta/2}$.  For the latter
law, its unnormalized mass after removing the terminal terms still
contains $\nu_{\varnothing}(1)\ge Z_*$; hence the same denominator
bound applies without an additional normalization factor.  Cycle support
and translation invariance in Proposition~\ref{tuple:laws} are exact
properties of these components and impose no further threshold.

For the leading component in Lemma~\ref{tuple:conditional-mixing}, the
incident graph has $m\le e_*$ edges.  Its list operator error is at most
$C_{\mathrm{gr}}p^{-\delta}$.  The split into $A\ge c_p^m/2$ and its
complement in that proof therefore bounds the squared conditional norm by
\[
 Z_0^{-1}C_{\mathrm{gr}}^2
 \bigl(2c_p^{-m}+4c_p^{-2m}\bigr)p^{-2\delta}\|f\|_2^2.
\]
Using $Z_0\ge Z_*$ and $c_p^{-m}\le(2k)^{e_*}$ bounds this by
\[
 C_{\mathrm{lead}}p^{-2\delta}\|f\|_2^2,
 \qquad
 C_{\mathrm{lead}}:=Z_*^{-1}C_{\mathrm{gr}}^2
 \bigl(2(2k)^{e_*}+4(2k)^{2e_*}\bigr).
\]
The case $m=0$ already has zero conditional norm.  For either mixture,
conditioning additionally on the component index and applying Jensen's
inequality bounds its squared conditional norm by
\[
 \bigl(C_{\mathrm{lead}}p^{-2\delta}
       +C_{\mathrm{mass}}p^{-\delta/2}\bigr)\|f\|_2^2
 \le (C_{\mathrm{lead}}+C_{\mathrm{mass}})
       p^{-\delta/2}\|f\|_2^2.
\]
Here the bound in every other component is $\|f\|_2^2$ because its
single-slot marginal is uniform.  As
$\delta/2-2\gamma=3\delta/10>0$, the desired squared bound
$p^{-2\gamma}\|f\|_2^2$ follows once
\[
 p^{\delta/2-2\gamma}\ge C_{\mathrm{lead}}+C_{\mathrm{mass}}.
\]

Choose an integer $P>P_1$ large enough that the displayed requirements
on $p$ above hold for every $p\ge P$, and that $P^{\gamma/2}\ge3$.
This is possible because $\delta>0$, $\alpha_{\mathrm{exc}}-4>0$, and
$\delta/2-2\gamma>0$, and all the constants depend only on $k,L,t$.
The product condition in \eqref{eq:9} can be
included in the same choice: for integer $P\ge2$,
\[
 \prod_{\substack{p\ge P\\p\ \mathrm{prime}}}(1+p^{-3})
 \le \exp\left(\sum_{n\ge P}n^{-3}\right)
 \le \exp\left(\frac{1}{2(P-1)^2}\right),
\]
whose last expression tends to one.  The covariance proof uses only the
exact progression identity and bijections when $p\nmid D$, so it imposes
no further restriction on $P$.  This proves the stated dependence of $P$.
Under Condition~\eqref{local:unit-condition}, Section~\ref{sec:local}
allows $P_1,L$ to be chosen from $k$ and each $E_p$ from $k,p$.  Also,
$t$ is chosen from $\delta$, which depends only on $k$.  The product defining $M_0$
then uses a degree-dependent set of primes and degree-dependent
exponents.  For a general fixed polynomial, $M_0$ retains any dependence
of the local exponents $E_p$ as well as that of $P_1$ and $L$.
\end{proof}

For the rest of the paper, fix a choice of $P$ supplied by
Proposition~\ref{tuple:threshold-dependence} and the corresponding $M_0$.
